% !TEX root = ../Dissertation.tex 

\chapter{{Zusammenfassung}} \label{Abschlussdiskussion}
Die vorgestellten empirischen Untersuchungen dienten in erster Linie dazu, grundlegende theoretische Annahmen über die Natur von deskriptiven und expressiven Intensivierern in einem breit gefächerten Rahmen zu überprüfen, um erstmals belastbare Aussagen über deren Zutreffen zu machen. Dazu wurde auf eine umfassende Kombination aus korpuslinguistischer und experimenteller Methodik zurückgegriffen. In diesem Kapitel werden die erzielten Ergebnisse mit Blick auf die zu testenden Hypothesen resümiert und final interpretiert. Dies geschieht in der Reihenfolge, in der sie in den einzelnen Kapiteln des empirischen Teils dieses Buchs beleuchtet wurden. 

Der erste Schritt der empirischen Untersuchung startete mit einer korpuslinguistischen Häufigkeitsanalyse. Im Kontext der in Kapitel \ref{Korpus} präsentierten Korpusstudie wurde eine umfangreiche Stichprobe von gegenwärtig verwendeten expressiven Intensivierern in Bezug auf ihre synchrone und diachrone Entwicklung hin untersucht und quantitativ eingeordnet. Dies geschah vor dem Hintergrund der folgenden beiden Hypothesen:

\begin{itemize}
\item[\textbf{\textsc{H1}}] Eine hochfrequente Verwendung von Intensivierern führt unweigerlich zum Verlust ihrer Expressivität.

\item[\textbf{\textsc{H2}}] Der Expressivitätsverlust geht mit einem Bedarf an neuen Ausdrücken einher, mit denen sich Sprecher auch weiterhin als expressiv, innovativ und originell darstellen können.
\end{itemize}

\noindent Während die Hypothese H2 in der Literatur auf \citet[274]{Baumgarten.1908} zurückgeht und annimmt, dass ständig neue Ausdrücke in die Intensiviererrolle überführt werden, leitet sich H1 aus einer Bemerkung von \citet[2]{Hauschild.1899} heraus ab. Dieser vertritt die Sichtweise, dass ein frequenter Gebrauch von Ausdrücken zulasten ihrer Expressivität geht, die in der Folge sukzessive schwindet. Dessen ungeachtet kommt für den Expressivitätsverlust neben der Frequenz auch ein zweiter Faktor infrage, nämlich die Zeit. Obwohl ihr möglicher Einfluss schon von \citet[59]{Tobler.1868} beschrieben wird, wurde die zeitliche Dimension meines Wissens in diesem Zusammenhang bislang nicht berücksichtigt. Um dem gerecht zu werden, wurde die Hypothese H1 anhand von zwei Faktoren spezifiziert:

\begin{itemize} 
\setlength{\itemindent}{1cm}
\item[\textbf{\textsc{H1-Freq}}] Eine hochfrequente Verwendung von Intensivierern führt zu einer pragmatischen Abnutzung und demzufolge sukzessive zum Verlust von Expressivität. 

\item[\textbf{\textsc{H1-Temp}}] Eine persistente Verwendung von Intensivierern führt zu einer pragmatischen Abnutzung und demzufolge sukzessive zum Verlust von Expressivität. 
\end{itemize}

Mit dem Ziel erste Hinweise auf die Gültigkeit der umrissenen Hypothesen zu erhalten, wurde zunächst eine Korpusanalyse durchgeführt, in der die Gebrauchsfrequenzen von 16 expressiven Intensivierern und dem deskriptiven Intensivierer \textit{sehr} ermittelt wurden. Als Untersuchungsgrundlage diente ein für diese Zwecke konzipiertes Textkorpus, das alle Ausgaben der Zeitschriften \textsc{Der Spiegel} und \textsc{Die Zeit} im Zeitraum von Januar 1950 (\textsc{Der Spiegel}) bzw. Januar 1953 (\textsc{Die Zeit}) bis Dezember 2017 umfasste. Das virtuelle Arbeitskorpus setzte sich in Summe aus 68 Jahrgängen, bestehend aus 133 Subkorpora mit 721\,609 Texten und ca. 562 Millionen Wörtern, zusammen.


Da mittels Korpusuntersuchung bloß das Vorkommen der Ausdrücke über die Zeit abgebildet werden kann, ohne dass Aussagen über deren Expressivitätsgrad gemacht werden, wurden die dadurch gewonnenen Ergebnisse als Spiegelbild des potenziellen Expressivitätsverlusts herangezogen. Im Hinblick auf die Hypothesen wurde davon ausgegangen, dass eine rückläufige diachrone Entwicklung ein hinreichendes Indiz für einen Expressivitätsverlust darstellt (vgl. H1), während das Auftreten neuer Intensivierer im Korpus einen Bedarf an neuen, unverbrauchten Ausdrücken signalisiert (vgl. H2).


Den Korpusfrequenzen zufolge lässt sich der Gebrauch des deskriptiven Intensivierers \textit{sehr} als eine Art Welle mit fallendem Beginn, leichter Plateau-Phase und ansteigendem Ende beschreiben (vgl. Abbildung \ref{Dia-sehr}). Quantitativ ist der Ausdruck mit insgesamt 20\,278,12 Treffern pro Million Textwörter mehr als zehn Mal so häufig in den Zeitschriften vertreten wie alle expressiven Intensivierer zusammen, die darin insgesamt 1\,719,36 Treffer pro Million Textwörter ausmachen. Weil es sich bei der eher konservativen Zeitungssprache dem Modell der Einteilung von Interaktionsformen von \citet[]{Koch.1996,Koch.2007} folgend um konzeptionell schriftlichen Sprachgebrauch handelt, der durch kommunikative Distanz gekennzeichnet ist, ist das starke Übergewicht von \textit{sehr} wenig überraschend. Die expressiven Intensivierer zeigen in Bezug auf ihre diachronen Häufigkeitsverteilungen einen beinahe linearen Anstieg im Korpus an: So hat sich ihr Gebrauch von Beginn des Untersuchungszeitraums bis zum Ende hin verdreifacht (vgl. Abbildung \ref{Dia-rest}). Dies wertete ich als einen Indikator für eine weniger stigmatisierte Verwendung der Ausdrücke in den Zeitschriften und als Signal stilistischer Veränderungen sowie zeitschrifteninterner Umstrukturierungen zugunsten mündlich orientierter Textsorten wie Interviews, Leserbriefen oder Kommentaren. Darüber hinaus weisen sie eine starke Variabilität sowohl in Bezug auf ihre Entwicklung (d.\,h. ansteigend, schwankend oder fallend) als auch auf ihre Häufigkeiten auf (vgl. Abbildung \ref{SyDia}). Dabei ist auch zu sehen, dass \textit{derb} der einzige expressive Intensivierer mit rückläufigem Entwicklungsverlauf ist. Aufgrund dessen bietet der Ausdruck als Einziger zweifelsfrei einen Grund zur Annahme des in der Hypothese H1 aufgegriffenen Expressivitätsverlusts \citep[vgl.][2]{Hauschild.1899}. Daneben wurde die Möglichkeit in Betracht gezogen, dass auch die kontinuierliche Frequenzzunahme der hochfrequenten Ausdrücke wie \textit{furchtbar} oder \textit{schrecklich} auf einen fortgeschrittenen Expressivitätsverlust hinweist und sich beim Expressivitätsverlust gegebenenfalls medienspezifische Unterschiede erkennen lassen. Diesbezüglich wurde gemutmaßt, dass ein Rückgang der Expressivität im schriftlichen Sprachgebrauch eventuell zu einer besseren Eignung, einer größeren Flexibilität im Gebrauch und einer höheren Verwendungshäufigkeit führt und die Motivation für die Verwendung von Intensivierern im Wesentlichen textsortenabhängig ist. Die im Korpus rezent aufgetretenen Intensivierer \textit{hammer}, \textit{krass} und \textit{mega} wurden außerdem als empirische Evidenz für die Gültigkeit der Hypothese H2 herangezogen, nach der permanent neue Ausdrücke zur Kompensation erforderlich sind \citep[vgl.][274]{Baumgarten.1908}. \par 
Um der Korrektheit der zur Diskussion stehenden Hypothesen H1-Temp und H1-Freq näher auf den Grund zu gehen, wurde der Untersuchungszeitraum in drei Intervalle à 20 Jahre unterteilt. In diesen wurden die expressiven Intensivierer unter Berücksichtigung der Faktoren Erstvorkommen und Frequenz im Detail inspiziert. Dazu wurden die Ausdrücke in jedem Zeitintervall in eine hochfrequente und eine niedrigfrequente Gruppe untergliedert. Unterstützt wurden die Einteilungen durch Dendrogramme, die mittels Clusteranalyse abgeleitet wurden. Im Zuge der Auswertung stellte sich heraus, dass sich die Mitglieder eines Clusters abgesehen von ähnlichen Frequenzen auch durch semantische Ähnlichkeiten auszeichnen. So setzt sich die Gruppe der hochfrequenten Intensivierer ergänzend zum frequentesten Intensivierer \textit{super} ausschließlich aus Ausdrücken mit klar negativer Semantik zusammen: \textit{unheimlich}, \textit{arg}, \textit{schrecklich}, \textit{furchtbar}, \textit{wahnsinnig} und \textit{verdammt}. Die hohen Frequenzen der Gruppenmitglieder wurden im Rahmen der Ergebnisinterpretation auf einen langen Gebrauch sowie einen entsprechend hohen Lexikalisierungsgrad zurückgeführt. Die Gruppe der niedrigfrequenten Ausdrücke beinhaltet hingegen neben den neuen Intensivierern \textit{hammer}, \textit{krass} und \textit{mega} zum einen \textit{derb} als rückläufigen Intensivierer sowie zum anderen \textit{ober}, \textit{fürchterlich}, \textit{scheiß}, \textit{voll} und \textit{sau}, die allesamt semantisch transparent sind. Zusammengefasst lassen die Ergebnisse der Korpusanalyse den Schluss zu, dass die Häufigkeiten der Ausdrücke zumindest partiell durch die Semantik bedingt zu sein scheinen: Während die semantische Transparenz den Gebrauch der Intensivierer im Zeitschriftenkorpus in einigen Fällen begünstigt hat (bspw. \textit{schrecklich}), hatte sie in anderen Fällen dagegen augenscheinlich eine eher hemmende Wirkung (bspw. \textit{scheiß}), was auf die Existenz eines semantischen Filters für die Schriftsprache schließen lässt.


Durch die Korpusstudie konnte also bereits ein Teil der Hypothesen bestätigt werden: So liefern die im Korpus neu erschienenen Intensivierer \textit{hammer}, \textit{krass} und \textit{mega} klare Evidenz für einen Bedarf an neuen Ausdrücken, mit denen Sprecher altgediente Formen ersetzen und weiterhin expressiv sein können. Dies steht im Einklang mit der Hypothese H2 nach \citet[274]{Baumgarten.1908}. Für die Hypothese H1 und den durch \citet[2]{Hauschild.1899} angeführten Expressivitätsverlust finden sich in Anbetracht des starken diachronen Anstiegs der expressiven Intensivierer weniger deutliche Anhaltspunkte. Lediglich \textit{derb} stellt in diesem Punkt eine Ausnahme dar, zeigt dieser Intensivierer doch als Einziger einen rückläufigen Entwicklungsverlauf auf, für den ein Expressivitätsverlust als garantiert ursächlich betrachtet wird. Wie es um die Expressivität der restlichen Ausdrücke steht, ließ sich mithilfe frequenzbasierter Korpusdaten jedoch nicht abschätzen, sodass weiterführende experimentelle Untersuchungen erforderlich waren.


Aus der Intention heraus Rückschlüsse auf die noch ungeklärten Hypothesen H1-Freq und H1-Temp ziehen zu können, wurde die quantitative Einordnung der Intensivierer anschließend um eine Einschätzung ihres Expressivitätsgrades erweitert. Dies geschah in Form einer experimentellen Untersuchung, die in Kapitel \ref{Experiment0} vorgestellt wurde und in der nach wie vor die Hypothesen H1-Freq und H1-Temp nach \citet[2]{Hauschild.1899} und \citet[59]{Tobler.1868} vordergründig waren. Ausgehend von den Befunden der Korpusanalyse ließen sich für das Experiment relevante Vermutungen ableiten, die es weitergehend zu überprüfen galt. So sind mutmaßlich sowohl im Korpus hochfrequente (z.\,B. \textit{super} oder \textit{verdammt}) als auch lange existierende Intensivierer (z.\,B. \textit{arg} und \textit{furchtbar}) grundsätzlich weniger expressiv als diejenigen Ausdrücke, die weniger frequent (z.\,B. \textit{sau} und \textit{voll}) oder neuer, d.\,h., erst im Laufe des Untersuchungszeitraums erstmals in intensivierender Funktion aufgetreten sind (z.\,B. \textit{hammer} oder \textit{krass}). Ein besonderes Augenmerk wurde auf den Intensivierer \textit{derb} gerichtet, der den Ergebnissen der Korpusstudie zufolge potenziell als nur wenig expressiv empfunden wird. Zusätzlich wurde im Rahmen des Experiments untersucht, ob die in der Literatur vorherrschenden Annahmen zutreffen, dass Expressivität gleichermaßen zu einem als höher wahrgenommenen Grad an Skalarität und Sprechereinstellung führt. Die dabei zu überprüfenden Hypothesen waren die folgenden:

\begin{itemize}
\item[\textbf{\textsc{H3}}] Expressive Intensivierer werden mit einem höheren Grad an Skalarität
assoziiert als deskriptive Intensivierer. 

\item[\textbf{\textsc{H4}}] Expressive Intensivierer geben eine Einstellungsbekundung des Sprechers
zu dem im Satz beschriebenen Sachverhalt wieder und werden in der Folge mit einem höheren Grad an Sprechereinstellung assoziiert als deskriptive Intensivierer. 
\end{itemize}

\noindent Der Hypothese H3 nach geben expressive Intensivierer einerseits einen höheren Grad einer Merkmalsausprägung an und verschieben den entsprechenden Skalenwert weiter nach rechts als deskriptive Intensivierer (vgl. \citealt[22]{Paradis.1997}; \citealt[133]{Gutzmann.2019a}). Andererseits dienen sie nach H4 dem Ausdruck der persönlichen Einstellung des Sprechers, weswegen sie bezüglich dieses Merkmals den Vermutungen nach ebenso mit einem höheren Grad assoziiert werden (vgl. \citealt[315]{Lang.1983}; \citealt[57]{Ortner.2014}).


Der mit den Ausdrücken einhergehende Grad an Skalarität und Sprechereinstellung wurde in einem Rating-Experiment erhoben, in dem eine Personenstichprobe kurze und einfach konstruierte Kopulasätze mit Intensivausdrücken, bestehend aus den berücksichtigten Intensivierern und Adjektiven, auf einer Skala bewertete. Die beiden Merkmale Skalarität und Sprechereinstellung wurden getrennt voneinander abgefragt, indem der Untersuchung zwei Skalen zugrunde gelegt wurden, auf denen die Versuchspersonen den jeweils zutreffenden Grad intuitiv einschätzten. Als Ergänzung zu den 16 expressiven Intensivierern wurde der deskriptive Intensivierer \textit{sehr} sowie die unmarkierte Variante des Satzes als Kontrollstufen miteingebunden. Damit sich das lexikalische Material nicht unerwünscht auf die Skaleneinschätzung auswirkt, wurden in einem mehrstufigen Verfahren semantisch leere Pseudowörter in Gestalt von Nomen und Adjektiven generiert. Dadurch konnte ein möglicher Einfluss verhindert und sichergestellt werden, dass die Beurteilung der Items ausschließlich durch die Intensivierer zustande kam. Weiterführend zur Hypothese H4 wurde die Überprüfung einer weiteren in diesem Zusammenhang relevanten Vermutung angestrebt:

\begin{itemize}
\item[\textbf{\textsc{H5}}] Die Einstellungsbekundung, die expressive Intensivierer zu dem im Satz beschriebenen Sachverhalt wiedergeben, ist evaluativer Natur und wird daher als positiv oder negativ wahrgenommen.
\end{itemize}

\noindent Diese Hypothese ergibt sich hauptsächlich aus einer Schlussfolgerung von \citet[135]{Gutzmann.2019a}, derzufolge sich die Einstellungsspiegelung expressiver Intensivierer aufgrund ihres evaluativen Charakters in Form einer positiven oder negativen Bewertung der bezeichneten Sachlage niederschlägt. Wie sich weiter ableiten lässt, werden deskriptive Intensivierer dahingehend mutmaßlich als überwiegend neutral wahrgenommen, drücken sie nun einmal nicht die persönliche Einstellung des Sprechers aus. Die Gültigkeit dieser Annahme wurde in einer dritten Aufgabe getestet, in der die Versuchspersonen die Polarität der Sprechereinstellung mithilfe einer komparativen Konstant-Summen-Skala evaluierten und eine zuvor festgelegte virtuelle Talersumme auf die drei Auswahlfelder `Positiv', `Neutral' und `Negativ' verteilten.


Die Ergebnisse des Experiments deuten schon auf den ersten Blick sowohl im Kontext der Skalarität als auch der Sprechereinstellung auf eine graduelle Anordnung der drei Stufen hin: expressive Intensivierer $\succ$ deskriptiver Intensivierer $\succ$ unmarkierte Variante (vgl. die Tabellen \ref{SkaWe1} und \ref{SkaWex}). Demnach werden Sätze mit expressivem Intensivierer in beiden Fällen mit einem bedeutend höheren Grad in Verbindung gebracht als Sätze mit deskriptivem Intensivierer oder unintensivierte Sätze. Besonders stark ausgeprägt ist der Unterschied bei der in Aufgabe II erfragten Sprechereinstellung. Hierbei gibt es keinen expressiven Intensivierer, dessen eingeschätzter Skalenwert unter dem des deskriptiven Intensivierers \textit{sehr} oder der unmarkierten Variante liegt. Dies stimmt mit der in der theoretischen Literatur vorherrschenden Auffassung überein, dass nur expressive Intensivierer eine Einstellungsbekundung des Sprechers zu dem im Satz beschriebenen Sachverhalt mitteilen \citep[vgl.][133]{Gutzmann.2019a}. Wie die Resultate in aller Deutlichkeit zeigen, sind die zu klärenden Hypothesen H3 und H4 nach \citet[22]{Paradis.1997} und \citet[133]{Gutzmann.2019a} sowie \citet[315]{Lang.1983} und \citet[57]{Ortner.2014} faktisch korrekt: Expressive Intensivierer werden mit einem signifikant höheren Grad an Skalarität und Sprechereinstellung assoziiert als deskriptive Intensivierer (vgl. die Tabellen \ref{LinReg0001} und \ref{LinReg0002}). Dies liefert Bestätigung für den angenommenen wahrnehmbaren Unterschied zwischen den Intensiviererklassen.

\largerpage
Außerdem wurde untersucht, ob zwischen der in Aufgabe I erfragten Skalarität und der in Aufgabe II fokussierten Sprechereinstellung ein korrelativer Zusammenhang besteht, wie \citet[133]{Gutzmann.2019a} in seiner Definition nahelegt. Das signifikante Ergebnis der Korrelationsanalyse spricht für eine positive Korrelation: Wurde der Grad an Sprechereinstellung von den Versuchspersonen auf der Skala hoch eingestuft, empfanden sie auch den Skalaritätsgrad als entsprechend höher. Somit sind die beiden Variablen \textsc{Skalarität} und \textsc{Sprechereinstellung} in dieser Untersuchung de facto klar voneinander abhängig. Ergänzend zur Korrelationsanalyse wurde ein lineares Regressionsmodell berechnet, mit dem der potenzielle Einfluss ebenso überprüft werden konnte. Dem Ergebnis der Analyse zufolge wirkt sich die Sprechereinstellung tatsächlich signifikant auf die eingeschätzte Skalarität aus (vgl. Tabelle \ref{LM01}), was mit der soeben erläuterten Vermutung in Einklang steht. Infolgedessen ist davon auszugehen, dass der Grad an Skalarität im Allgemeinen umso höher evaluiert wird, je stärker der Grad an Sprechereinstellung wahrgenommen wird.


Die Ergebnisse zu Aufgabe III, in der es um die Beurteilung der Einstellungspolarität ging, lassen den Schluss zu, dass sowohl ein hoher Grad an Skalarität als auch ein hoher Grad an Sprechereinstellung in einer signifikanten Abweichung von einer als neutral wahrgenommenen Sprechereinstellung resultiert (vgl. Tabelle \ref{LinReg0004}): Je höher beide Merkmale also eingeschätzt wurden, desto stärker bewerteten die Versuchspersonen die Einstellungspolarität als positiv oder negativ. Bei der Einzelbetrachtung der Ausdrücke fiel ferner auf, dass die eingeschätzte Sprechereinstellung vieler expressiver Intensivierer tendenziell negativ ausfällt, was zu einer Überrepräsentatiozn negativer Polarität führt (vgl. Tabelle \ref{ValenzInten}). In diesem Zusammenhang wurde gemutmaßt, dass dies möglicherweise auf die oftmals offenkundige negative Semantik der expressiven Ausdrücke zurückzuführen ist, die dem Anschein nach die Beurteilung beeinflusst. Gestützt wird diese Einschätzung durch die als positiv bewerteten Ausdrücke. Dabei handelt es sich nahezu ausnahmslos um solche, die als Interjektionen bzw. innerhalb von Ausrufesätzen in prädikativer Verwendung auftreten können und in dieser Stellung ausschließlich positiv konnotiert sind (bspw. \textit{Hammer!} oder \textit{Das ist ja super!}). Im Gegensatz dazu wurden der deskriptive Intensivierer \textit{sehr} und die unmarkierte Variante des Satzes vermehrt als neutral beurteilt (vgl. Abbildung \ref{Aufg3.1}). Dies untermauert die zu testende Hypothese H5 nach \citet[135]{Gutzmann.2019a} und demonstriert weiterhin, dass eine komparative Konstant-Summen-Skala wie die in dieser Untersuchung verwendete ein geeignetes Werkzeug ist, um die Polarität der Sprechereinstellung experimentell zu erheben.


\largerpage
Ergänzend zu den Merkmalen Skalarität, Sprechereinstellung und Einstellungspolarität wurde in der Untersuchung der Rolle von Frequenz nachgegangen. Um Hinweise auf das Zutreffen der Hypothese H1-Freq nach \citet[2]{Hauschild.1899} zu erhalten, wurden die expressiven Intensivierer gemäß der Korpusfrequenzen des letzten Zeitabschnitts (d.\,h. dem Zeitraum von 1990 bis 2017) in eine niedrig- und eine hochfrequente Gruppe eingeteilt (vgl. Tabelle \ref{Fre3.1}), deren Unterscheidung bei der Regressionsanalyse als Ausgangspunkt gewählt wurde. Bei der in Aufgabe I untersuchten Skalarität zeigt sich ein signifikanter Unterschied zwischen den beiden Frequenzgruppen (vgl. Tabelle \ref{LinRea0100}). Dem Regressionsmodell nach werden die niedrigfrequenten Intensivierer mit einem deutlich niedrigeren Skalaritätsgrad assoziiert als die hochfrequenten Intensivierer. Dieser Befund lässt den Schluss zu, dass die wahrnehmbare Skalarität eines Ausdrucks mit steigender Frequenz zunimmt. Auch die Ergebnisse der in Aufgabe II ermittelten Sprechereinstellung sprechen für einen signifikanten Unterschied zwischen hoch- und niedrigfrequenten Intensivierern (vgl. Tabelle \ref{LinRegzaz0}). Demzufolge wird der Grad an Sprechereinstellung bei den niedrigfrequenten Intensivierern als wesentlich höher empfunden als bei den hochfrequenten Intensivierern. Je frequenter und erwartbarer ein Intensivierer also war, desto niedriger wurde der durch ihn ausgedrückte Grad an Sprechereinstellung von den Versuchspersonen wahrgenommen. Dieser Befund befürwortet die Annahme, dass eine hohe Frequenz auf Kosten der Expressivität von Ausdrücken geht, was die zu überprüfende Hypothese H1-Freq nach \citet[2]{Hauschild.1899} erhärtet.


Analog zum Vorgehen bei der Frequenz wurden die expressiven Intensivierer auch bezüglich der Persistenz in zwei Gruppen eingeordnet. Dies geschah auf Grundlage ihres Erstvorkommens im Korpus, was bedeutet, dass die drei rezent in Erscheinung getretenen Intensivierer \textit{hammer}, \textit{krass} und \textit{mega} bei der Datenmodellierung von den restlichen Ausdrücken separiert wurden, um Rückschlüsse auf die Gültigkeit der Hypothese H1-Temp nach \citet[59]{Tobler.1868} ziehen zu können. Die Ergebnisse zur in Aufgabe I ermittelten Skalarität sind signifikant (vgl. Tabelle \ref{LinRea0100}). Demnach wird der Skalaritätsgrad der neuen Intensivierer auf der Skala bedeutend höher verortet als der der älteren Intensivierer. Dagegen tritt in Bezug auf die in Aufgabe II untersuchte Sprechereinstellung kein signifikanter Unterschied zwischen alten und neuen Ausdrücken auf (vgl. Tabelle \ref{LinRegzaz0}), weswegen die Hypothese {H1-Temp} aufgrund fehlender empirischer Evidenz verworfen werden muss. Wie in diesem Zusammenhang erläutert, ist dieser Befund unter Beachtung des Ungleichgewichts, das zwischen den Persistenzgruppen herrscht (drei neue vs. 13 ältere Intensivierer), bloß eingeschränkt aufschlussreich. So ist zu erwägen, dass sich dieses auf die inferenzstatistischen Auswertungen ausgewirkt und den erwarteten Effekt überlagert hat. Aus diesem Grund besteht die Möglichkeit, dass das Experiment mit einem ausbalancierten Verhältnis aus alten und neuen Ausdrücken durchaus zu einem signifikanten Unterschied zwischen den Persistenzgruppen geführt hätte, auch wenn über die Korrektheit dieser Mutmaßung an dieser Stelle keine fundierte Aussagen gemacht werden können. Insofern sollte in nachfolgenden Studien ein ausgewogeneres Verhältnis aus alten und neuen Intensivierern angestrebt werden, um weitere Anhaltspunkte für die nach wie vor ungeklärte Hypothese H1-Temp nach \citet[59]{Tobler.1868} zu erhalten.

In einer weiteren Untersuchung wurde der Annahme eines wahrnehmbaren Unterschiedes zwischen deskriptiven und expressiven Intensivierern noch einmal gesondert nachgegangen. Darin wurde eine weiterführende Hypothese getestet, die sich unmittelbar aus diesem Unterschied ableiten lässt: 

\begin{itemize}
\item[\textbf{\textsc{H6}}] Innerhalb der Intensiviererklassen lassen sich sowohl im Hinblick auf Skalarität als auch auf Sprechereinstellung graduelle Abstufungen zwischen den Ausdrücken beobachten, die semantisch bedingt sind. 
\end{itemize}

\noindent Aus dieser Überlegung heraus folgt, dass es nicht nur zwischen, sondern auch innerhalb der Intensiviererklassen Unterschiede bzw. graduelle Abstufungen des durch die Intensivierer ausgedrückten Intensitätsgrades gibt, die u.\,a. \citet[87]{Biedermann.1969} und \citet[408]{Breindl.2007} besprechen. In diesem Sinne lassen sich die Ausdrücke mutmaßlich in Abhängigkeit ihrer Semantik auf einer kontinuierlichen, längs organisierten Skala anordnen. Dabei besteht die Möglichkeit, dass die Expressivität Einfluss auf den Skalaritätsgrad nimmt, der in der Folge als stärker wahrgenommen wird \citep[vgl.][133]{Gutzmann.2019a}.


Um die Hypothese H6 zu testen, wurde eine Stichprobe aus jeweils vier deskriptiven und expressiven Intensivierern mit unterschiedlich wahrnehmbarer Intensität evaluiert. Dies geschah in zwei Rating-Experimenten mit parallelem Untersuchungsaufbau, in denen eine Personenstichprobe Sätze mit deskriptiven und expressiven Intensivausdrücken hinsichtlich Skalarität (Experiment 2a) und Sprechereinstellung (Experiment 2b) bewertete. Durch die Parallelität konnten die Ergebnisse der beiden Experimente miteinander verglichen und auf einen möglichen korrelativen Zusammenhang zwischen Skalarität und Sprechereinstellung hin überprüft werden.


Die Untersuchungsergebnisse bestätigen ebenfalls die oben angeführten Hypothesen H3 und H4, wurden Sätze mit expressiven Intensivierern doch sowohl im Kontext der Skalarität als auch der Sprechereinstellung von den Versuchspersonen auf der Skala signifikant höher positioniert als Sätze mit deskriptiven Intensivierern (vgl. die Tabellen \ref{LinReg01} und \ref{LinReg2}). Kurz gefasst sprechen auch diese Resultate in Summe ganz klar für die Annahmen von \citet[22]{Paradis.1997} und \citet[133]{Gutzmann.2019a} sowie \citet[315]{Lang.1983} und \citet[57]{Ortner.2014} und folglich die Gültigkeit der Hypothesen H3 und H4.


%Wenngleich der Unterschied bei der in \mbox{Experiment 2a} abgefragten Skalarität nicht übermäßig stark ausgeprägt war (vgl. \mbox{Tabelle \ref{tab:Skalenwert1}}), war er bei der in Experiment 2b fokussierten Sprechereinstellung dagegen umso größer (vgl. \mbox{Tabelle \ref{SkaWe_E2}}).


Weil der Fokus der Untersuchung auf den potenziellen klasseninternen Abstufungen lag, wurden die Intensiviererklassen in einem weiteren Analyseschritt aufgeschlüsselt und ihre Mitglieder im Detail inspiziert. Im Zuge dessen zeigten sich bei den deskriptiven Intensivierern die vermuteten graduellen Abstufungen: Die Ausdrücke ließen sich in Abhängigkeit ihrer Semantik hinsichtlich beider Merkmale als schwach, mittel oder stark klassifizieren (vgl. Tabelle \ref{deIn}). Dieser Befund weist auf den ersten Blick auf ein Zutreffen der Hypothese H6 hin, die von solchen klasseninternen Unterschieden ausgeht. Dass sich die Anordnung der Ausdrücke je nach untersuchtem Merkmal, d.\,h. Skalarität oder Sprechereinstellung, unterscheidet, lässt außerdem vermuten, dass die Versuchspersonen sie als verschieden wahrgenommen und dementsprechend unterschiedlich bewertet haben. Überraschenderweise wurde der Grad an Sprechereinstellung bei den deskriptiven Intensivierern auf der Skala nicht so niedrig beurteilt, wie aufgrund der fehlenden Expressivität zu erwarten war. Dies vermittelt den Eindruck, dass deskriptive Intensivierer zu einem gewissen Grad ebenso als einstellungsausdrückend wahrgenommen werden, gleichwohl die Expressivität bei ihnen nicht als Erklärung für die vergleichsweise hohe Skaleneinschätzung infrage kommen kann. Auch bei den expressiven Intensivierern sind graduelle Unterschiede auszumachen, in deren Folge die Ausdrücke in Abhängigkeit ihres Skalaritätsgrades als schwach, mittel oder stark klassifiziert werden können (vgl. Tabelle \ref{exIn}). Im Kontext der Sprechereinstellung lassen sie sich hingegen zu zwei Stufen mit je zwei Mitgliedern zusammenfassen: stark oder schwach einstellungsausdrückend. Dieses Ergebnis liefert die erste empirische Evidenz für die von \citet[87]{Biedermann.1969} und \citet[408]{Breindl.2007} angenommenen klasseninternen Abstufungen und bekräftigt die zur Diskussion stehende Hypothese H6. Ungeachtet der Ergebnisse wurde darauf hingewiesen, dass im Rahmen einer solchen experimentellen Erhebung keine trennscharfe oder isolierte Betrachtung beider Bedeutungsdimensionen erfolgen kann und davon auszugehen ist, dass die Expressivität trotz zielgerichteter Aufgabenstellung auf die Einschätzung des Skalaritätsgrades einwirkt. Demnach ist es bei einer Untersuchung von Skalarität aus methodischen Gründen nicht möglich, die Merkmalsausprägung eines im Satz beschriebenen Gegenstands ohne Einfluss der expressiven Komponente des jeweiligen Intensivierers zu testen und den Expressivitätsgrad gänzlich vom Skalaritätsgrad zu trennen. Deswegen bleibt das genaue Ausmaß der Auswirkung von Expressivität auf die Einschätzung von Skalarität über die Untersuchung hinaus ungeklärt.


Zusätzlich wurde getestet, ob die beiden Variablen \textsc{Skalarität} und \textsc{Sprechereinstellung} in einer Korrelation stehen. Dem Ergebnis der Korrelationsanalyse nach war dies nicht der Fall. Dieses Resultat ist insofern unerwartet, als gemäß der Definition von Expressivität nach \citet[133]{Gutzmann.2019a} doch von einem korrelativen Zusammenhang auszugehen ist. Dieser nimmt an, dass ein erhöhter Grad an Sprechereinstellung zugleich von einem höheren Grad an Skalarität begleitet wird. Auch das ergänzend berechnete lineare Regressionsmodell deutete auf keinen signifikanten Einfluss von Sprechereinstellung auf Skalarität hin (vgl. Tabelle \ref{LM02}), was ebenfalls keinen konkreten Anhaltspunkt für die Vermutung im Sinne von \citet[133]{Gutzmann.2019a} liefert. Dass die angenommene Korrelation durch die Untersuchung nicht nachgewiesen werden kann, hat aber bloß eine begrenzte Aussagekraft, wurden in die Analyse doch lediglich die Mittelwerte der expressiven Intensivierer miteinbezogen. Dagegen wurden die Werte der deskriptiven Intensivierer aufgrund fehlender Expressivität zur Gänze außen vor gelassen. Vor diesem Hintergrund wurde jedoch klargestellt, dass eine Korrelationsanalyse mit lediglich vier Datenpunkten allenfalls eingeschränkt interpretierbar ist und bei einer größeren Datenmenge möglicherweise doch eine signifikante Korrelation hätte festgestellt werden können. Eine solche signifikante Korrelation zeigte sich bspw. bei dem in Kapitel \ref{Experiment0} besprochenen Experiment, in dessen Korrelationsanalyse die Werte von 16 expressiven Intensivierern einbezogen wurden: Deren Ergebnis spricht durchaus für einen signifikanten Zusammenhang zwischen den Merkmalen Skalarität und Sprechereinstellung. In der Folge bleibt festzuhalten, dass ihre weitere Untersuchung nachfolgenden Experimenten vorbehalten bleiben muss.


Zusammengefasst bestätigen die in diesem Buch vorgestellten Ergebnisse nahezu alle zur Diskussion stehenden Hypothesen: Wird die Hypothese H1-Freq nach \citet[2]{Hauschild.1899} durch einen signifikanten Unterschied zwischen den Frequenzgruppen gestützt, untermalen die im Korpus neu erschienenen Intensivierer \textit{hammer}, \textit{krass} und \textit{mega} die Korrektheit der Hypothese H2 nach \citet[274]{Baumgarten.1908}, derzufolge durch den Expressivitätsverlust alter Intensivierer fortwährend neue Ausdrücke zur Kompensation erforderlich sind. Darüber hinaus findet sich ebenso empirische Evidenz für die Gültigkeit der Hypothesen H3 und H4 nach \citet[22]{Paradis.1997} und \citet[133]{Gutzmann.2019a} sowie \citet[315]{Lang.1983} und \citet[57]{Ortner.2014}, nach denen Expressivität gleichermaßen mit einem höheren Grad an Skalarität und Sprechereinstellung einhergeht. Abgesehen davon kann auch die Hypothese H5 für gültig befunden werden, die sich aus \citet[135]{Gutzmann.2019a} ableitet. Dieser postuliert, dass die Einstellungsbekundung, die expressive Intensivierer kommunizieren, evaluativ sei und sich in Form einer positiven oder negativen Bewertung des beschriebenen Sachverhalts niederschlage. Weiterhin gibt es stichhaltige Hinweise auf das Zutreffen der Hypothese H6, die auf Basis der theoretischen Literatur (vgl. bspw. \citealt[87]{Biedermann.1969}; \citealt[408]{Breindl.2007}) von klasseninternen Unterschieden ausgeht. Ebensolche graduellen Abstufungen  treten in der Untersuchung sowohl im Kontext der Skalarität als auch der Sprechereinstellung auf. Folglich konnte einzig für die Hypothese H1-Temp nach \citet[59]{Tobler.1868}, laut der langer Gebrauch nachteilig für die Expressivität von Ausdrücken ist, keine harte empirische Evidenz gefunden werden. So weist die Regressionsanalyse der in Kapitel \ref{Experiment0} beschriebenen Untersuchung in Bezug auf die Sprechereinstellung keinen signifikanten Unterschied zwischen den Persistenzgruppen auf. In Anbetracht dessen muss davon ausgegangen werden, dass die drei im Korpus rezent aufgetretenen Intensivierer \textit{hammer}, \textit{krass} und \textit{mega} von den Versuchspersonen als nicht signifikant expressiver empfunden wurden als die älteren Intensivierer (n = 13). Dieses Resultat ist unter Berücksichtigung der niedrigen Anzahl von neuen Ausdrücken allerdings nur eingeschränkt aufschlussreich, hat das daraus resultierende Ungleichgewicht im Umfang der Persistenzgruppen doch vermutlich auf die inferenzstatistischen Berechnungen ein- und dem erwarteten Effekt entgegengewirkt. Demnach ist es vorstellbar, dass die Ergebnisse mit einem ausgeglichenen Verhältnis aus alten und neuen Intensivierern in diesem Punkt
anders ausfallen würden. Dieser methodische Missstand könnte in Folgeexperimenten beachtet werden. Auch kann die fehlende Korrelation bei der in Kapitel \ref{Experimente} präsentierten Untersuchung bloß mit Vorsicht interpretiert werden. Wie oben beschrieben, bedarf es diesbezüglich ebenfalls einer größeren Datenmenge als derjenigen, die in diese Analyse einbezogen wurde. Auch wenn sich in diesen beiden Experimenten kein signifikanter Zusammenhang zwischen den Variablen \textsc{Skalarität} und \textsc{Sprechereinstellung} findet, bedeutet dies nicht, dass die beiden Merkmale grundsätzlich nicht korrelieren, wie das Ergebnis der in Kapitel \ref{Experiment0} dargelegten Korrelationsanalyse nahelegt.


%Abschließend gilt es, die diesem Buch übergeordnete, titelgebende Frage zu beantworten, die auf eine Einschätzung von \citet[133]{Gutzmann.2019a} zurückgeht: Ist \textit{sau cool} cooler als \textit{sehr cool}? Wie die Ergebnisse der experimentellen Untersuchungen mit deutlicher statistischer Signifikanz zeigen, scheint diese in der Tat zutreffend zu sein. In diesem Sinne kann schlussendlich ganz klar festgehalten werden: \textit{Sau cool} ist zweifellos cooler als \textit{sehr cool}.





%Bei der anschließend vorgenommenen Einzelbetrachtung der Intensivierer fiel im Kontext von Experiment 2a eine vorherrschende Trennung nach deskriptiv und expressiv auf (vgl. \mbox{Tabelle \ref{Skalenwert3}}). Dennoch war in beiden Intensivierklassen eine Ausnahme auszumachen, nämlich \textit{überaus} aufseiten der deskriptiven sowie \textit{voll} aufseiten der expressiven Intensivierer (vgl. \mbox{Abbildung \ref{Bar1}}). Während ersterem mit durchschnittlich 62,36\,\% ein verhältnismäßig hoher Skalaritätsgrad zugesprochen wurde, wurde \textit{voll} mit durchschnittlich 59,69\,\% von den Versuchspersonen dahingehend auf der Skala merklich niedrig verortet. Für diesen Befund wurden zwei Interpretationsmöglichkeiten der Hypothese H4 offengelegt, von denen die gemäßigtere eine mögliche Erklärung für dieses scheinbar abweichende Verhalten der beiden Ausdrücke bot (vgl. \mbox{Tabelle \ref{schwach}}). Die Einzelbetrachtung der Intensivierer von Experiment 2b zeigte hingegen keine solchen vermeintlichen Ausreißer, was bedeutet, dass sich in Bezug auf Sprechereinstellung erwartungsgemäß eine markante Trennung nach Intensiviererklasse herauskristallisierte (Tabelle \ref{Einzeln_2} und \mbox{Abbildung \ref{Bar2}}).



%in Übereinstimmung mit den Erwartungen von den Versuchspersonen auf der Skala grundsätzlich höher verortet als Sätze mit deskriptivem Intensivierer oder unintensivierte Sätze. Bei der Einzelbetrachtung von Aufgabe I, die die Skalarität erfragte, fiel auf, dass sich der deskriptive Intensivierer \textit{sehr} prozentual kaum von manchen expressiven Intensivierern unterschied (bspw. \textit{arg} (59,87\,\%) oder \mbox{\textit{derb} (63,36\,\%))}. Ferner war das Skalenspektrum der expressiven Intensivierer in dem Sinne eng gefasst, dass alle Ausdrücke den Einschätzungen der Versuchspersonen nach in einem durchschnittlichen Skalenbereich von \mbox{59,89\,\% (\textit{arg})} und 74,66\,\% (\textit{mega}) angesiedelt wurden (vgl. Tabelle \ref{SkaWe2} und Abbildung \ref{Skalar7}). Das zu \mbox{Aufgabe I} berechnete Regressionsmodell zeigte signifikante Haupteffekte für die Prädiktoren \textsc{Bedingung} und \textsc{Position} an (vgl. Tabelle \ref{LinReg001}). So war zum einen der Unterschied zwischen Sätzen mit expressivem Intensivierer und unintensivierten Sätzen signifikant (\mbox{$\chi$$^{2}$\,(1) = 54,42}, \mbox{\textit{p} < 0,001}), ebenso wie zum anderen auch der Unterschied zwischen Sätzen mit expressivem Intensivierer und Sätzen mit deskriptivem Intensivierer signifikant war \mbox{($\chi$$^{2}$\,(1) = 3,74}, \textit{p} < 0,05). Der signifikante Positionseffekt deutete zudem auf eine Abnahme der Skalarität in Abhängigkeit der Position eines Items im Experiment an \mbox{($\chi$$^{2}$\,(1) = 6,90}, \mbox{\textit{p} < 0,01}). Ähnliches zeigte sich beim finalen Modell, bei dem 'Unmarkiert' als Referenzlevel gesetzt wurde (vgl. Tabelle \ref{LinReg010}): Dieses enthielt ebenfalls signifikante Haupteffekte für die Prädiktoren \textsc{Bedingung} und \textsc{Position}. Somit wurden die unintensivierten Sätze nicht nur signifikant niedriger evaluiert als mittels expressivem Intensivierer modifizierte Sätze ($\chi$$^{2}$\,(1) = 54,42, \textit{p} < 0,001), sondern auch als mittels deskriptivem Intensivierer modifizierte Sätze \mbox{($\chi$$^{2}$\,(1) = 44,08}, \mbox{\textit{p} < 0,001}). Daneben gab es erneut einen signifikanten Positionseffekt \mbox{($\chi$$^{2}$\,(1) = 6,89}, \mbox{\textit{p} < 0,01}).


%Im Kontext der in Aufgabe II fokussierten Sprechereinstellung war der Unterschied zwischen dem deskriptiven Intensivierer \textit{sehr} und den expressiven Intensivierern deutlich stärker ausgeprägt (vgl. Tabelle \ref{SkaWex} und \mbox{Abbildung \ref{Expressiv7}}). So gab es an dieser Stelle keinen expressiven Intensivierer, dem von den Versuchspersonen eine niedrigere Skalenposition zugeschrieben wurde als dem deskriptiven Intensivierer. Auch bei dieser Aufgabe war das Skalenspektrum eng gefasst, lag der am niedrigsten bewertete expressive Intensivierer \textit{arg} bei durchschnittlich 56,94\,\%, während der auf der Skala am höchsten verortete expressive Intensivierer \textit{scheiß} im Schnitt bei 78,12\,\% lag. Das zu \mbox{Aufgabe II} berechnete Regressionsmodell wies signifikante Haupteffekte für den Prädiktor \textsc{Bedingung} auf (vgl. Tabelle \ref{LinReg002}), was bedeutet, dass sich Sätze mit expressivem Intensivierer in Bezug auf Sprechereinstellung signifikant von Sätzen mit deskriptivem Intensivierer ($\chi$$^{2}$\,(1) = 36,86, \textit{p} < 0,001) und unintensivierten Sätzen ($\chi$$^{2}$\,(1) = 51,58, \textit{p} < 0,001) unterschieden. Wiederholt wurde das Referenzlevel ergänzend auf 'Unmarkiert' gesetzt, wobei auch hier signifikante Haupteffekte für den Prädiktor \textsc{Bedingung} zu verzeichnen waren (vgl. Tabelle \ref{LinReg020}). So war nicht nur der Unterschied zwischen unintensivierten Sätzen und Sätzen mit expressivem Intensivierer signifikant \mbox{($\chi$$^{2}$\,(1) = 51,58}, \mbox{\textit{p} < 0,001}), sondern auch der Unterschied zwischen unintensivierten Sätzen und Sätzen mit deskriptivem Intensivierer ($\chi$$^{2}$\,(1) = 13,08, \textit{p} < 0,001).


%Außerdem wurde untersucht, ob zwischen der in Aufgabe I erfragten Skalarität und der in Aufgabe II fokussierten Sprechereinstellung ein korrelativer Zusammenhang bestand. Zur Quantifizierung des potenziellen linearen Merkmalszusammenhangs wurde der Korrelationskoeffizient \textit{r} nach Pearson berechnet, dessen \textit{Estimate} auf einen Zusammenhang hindeutete: \mbox{|r| = 0,83}. Das Ergebnis des zugehörigen \textit{t}-Tests war signifikant (\mbox{\textit{t}\,(16) = 5,94}, \textit{p} < 0,001), was eine starke positive Korrelation zwischen Skalarität und Sprechereinstellung nahelegte: Wenn der Skalaritätsgrad auf der Skala hoch eingestuft wurde, wurde auch der Grad an Sprechereinstellung als hoch empfunden. Dieser Befund stimmte mit den Erwartungen überein, die im Vorhinein der Untersuchung an die expressiven Intensivierer gestellt wurden.


%Im nächsten Schritt wurden die expressiven Intensivierer gemäß der Korpusfrequenzen des letzten Zeitabschnitts (d.\,h. dem Zeitraum von 1990 bis 2017) in eine niedrig- und eine hochfrequente Gruppe eingeteilt (vgl. Tabelle \ref{Fre3.1}). Diese Gruppen wurden im Regressionsmodell als Ausgangspunkt gewählt, um der Gültigkeit der Hypothese H1-Freq nachzugehen. Das in Tabelle \ref{LinRea} gelistete finale Modell zu Aufgabe I enthielt einen Haupteffekt für den Prädiktor \textsc{Position} ($\chi$$^{2}$\,(1) = 5,95, \textit{p} < 0,05). Für die Frequenz war an dieser Stelle allerdings keinen Effekt auszumachen, was bedeutet, dass sie nachweislich keinen Einfluss auf den Skalaritätsgrad hatte. Entsprechend fand sich auch kein korrelativer Zusammenhang zwischen Frequenz und Skalarität. So wies der \textit{Estimate} des Korrelationskoeffizienten \textit{r} bereits auf Unkorreliertheit hin: \mbox{|r| = -0,19}. Das Ergebnis des zugehörigen \textit{t}-Tests war ebenfalls nicht signifikant (\mbox{\textit{t}\,(15) = -0,74}, \textit{p} > 0,05). Anders sah dies dagegen im Kontext der in Aufgabe II erfragten Sprechereinstellung aus. Das dafür berechnete finale Regressionsmodell enthielt einen signifikanten Haupteffekt für den Prädiktor \textsc{Frequenzgruppe} (vgl. \mbox{Tabelle \ref{LinRegzaz}}). Demnach wurde der Grad an Sprechereinstellung bei den niedrigfrequenten Intensivierern als signifikant höher empfunden als bei den hochfrequenten Intensivierern ($\chi$$^{2}$\,(1) = 3,35, \mbox{\textit{p} < 0,05)}. Des Weiteren trat eine negative Korrelation zwischen Frequenz und Sprechereinstellung auf, wie der \textit{Estimate} des Korrelationskoeffizienten \textit{r} nahelegte: \mbox{|r| = -0,63}. Das Ergebnis des \textit{t}-Tests war signifikant \mbox{(\textit{t}\,(15) = -3,17}, \textit{p} < 0,01). Je frequenter und erwartbarer ein Intensivierer also war, desto niedriger wurde der durch ihn ausgedrückte Grad an Sprechereinstellung wahrgenommen. Dies lieferte die erste empirische Evidenz für das Zutreffen der Hypothese H1-Freq nach \mbox{\citet[2]{Hauschild.1899}}.


%Analog zum Vorgehen bei Frequenz wurden die expressiven Intensivierer auch im Kontext von Persistenz in zwei Gruppen eingeteilt. Dies geschah auf Grundlage ihres Erstvorkommens im Korpus, was bedeutet, dass die drei rezent in Erscheinung getretenen Intensivierer \textit{hammer}, \textit{krass} und \textit{mega} bei der Datenmodellierung von den restlichen Intensivierern separiert wurden, um Rückschlüsse auf das Zutreffen von H1-Temp ziehen zu können. Das finale Regressionsmodell zu Aufgabe I enthielt signifikante Haupteffekte für die Prädiktoren \textsc{Persistenzgruppe} und \textsc{Position} (vgl. Tabelle \ref{LinRegzy}). So wurde zum einen der Skalaritätsgrad der neuen Intensivierer auf der Skala signifikant höher verortet als der der älteren Intensivierer (\mbox{$\chi$$^{2}$\,(1) = 13,84}, \mbox{\textit{p} < 0,001)}, während der signifikante Positionseffekt zum anderen anzeigte, dass die wahrgenommene Skalarität mit fortschreitendem Experiment abgenommen hat \mbox{($\chi$$^{2}$\,(1) = 5,14}, \mbox{\textit{p} < 0,05)}. Das finale Regressionsmodell, das für die in \mbox{Aufgabe II} erfragte Sprechereinstellung berechnet wurde, beinhaltete keine signifikanten Haupteffekte (vgl. Tabelle \ref{LinRegzaz2}). Somit wirkte sich die Persistenz nachweislich nicht auf den Grad an Sprechereinstellung aus, weswegen die Gültigkeit der Hypothese H1-Temp nach \citet[59]{Tobler.1868} aufgrund fehlender empirischer Evidenz an dieser Stelle verworfen wurde.


%Die Ergebnisse zu Aufgabe III, in der es um die Beurteilung der Einstellungspolarität ging, zeigten, dass sich Sätze mit expressivem Intensivierer insofern von Sätzen mit deskriptivem Intensivierer oder unintensivierten Sätzen abhoben, als deren Beurteilung verstärkt in Richtung 'negativ' tendierte. Dagegen wurden die beiden anderen Bedingungen von den Versuchspersonen vielmehr als neutral eingeschätzt (vgl. Tabelle \ref{Bedin}). Gestützt wurde diese erste Einschätzung durch die Einzelbetrachtung der Ausdrücke (vgl. \mbox{Tabelle \ref{ValenzInten}} und \mbox{Abbildung \ref{Aufg3.1}}). Dabei stellte sich heraus, dass die Interpretation der Intensivierer mit klar negativem Ursprung (bspw. \textit{furchtbar} oder \textit{schrecklich}) stark in die negative Richtung tendierte ($\hat{=}$ 44,44\,\%), wohingegen positiv konnotierte Ausdrücke (bspw. \textit{mega} oder \textit{super}) oder solche, die im Deutschen interjektionell oder innerhalb von Ausrufesätzen in prädikativer Verwendung auftreten können (bspw. \textit{hammer} oder \textit{wahnsinnig}), überwiegend als positiv beurteilt wurden ($\hat{=}$ 27,78\,\%). Eine Tendenz in Richtung 'neutral' gab es in zwei Fällen, nämlich beim deskriptiven Intensivierer \textit{sehr} sowie der unmarkierten Variante ($\hat{=}$ 11,11\,\%). Drei expressive Intensivierer wiesen darüber hinaus ein ausgewogenes Verhältnis auf, bei dem keine ausgeprägte Tendenz zugunsten einer Richtung auszumachen war: \textit{krass}, \textit{unheimlich} und \textit{voll} \mbox{($\hat{=}$ 16,67\,\%)}. Das zu Aufgabe III berechnete finale Modell enthielt signifikante Haupteffekte für die Prädiktoren \textsc{Expressivität}, \textsc{Skalarität} und \textsc{Position} (vgl. Tabelle \ref{LinReg004}). Dementsprechend wurden Intensivierer, die mit einem hohen Grad an Sprechereinstellung assoziiert wurden, als signifikant positiver oder negativer evaluiert (\mbox{$\chi$$^{2}$\,(1) = 39,67}, \mbox{\textit{p} < 0,001)}, ebenso wie auch diejenigen Intensivierer als signifikant positiver oder negativer bewertet wurden, denen ein hoher Skalaritätsgrad zugeschrieben wurde (\mbox{$\chi$$^{2}$\,(1) = 9,25}, \mbox{\textit{p} < 0,01)}. Dem signifikanten Positionseffekt nach wurde die Abweichung von 'neutral' umso stärker empfunden, je weiter fortgeschritten die Position eines Items im Experiment war (\mbox{$\chi$$^{2}$\,(1) = 4,11}, \mbox{\textit{p} < 0,05)}. Insgesamt sprachen die Befunde von Aufgabe III für die Gültigkeit der Hypothese H3, nach der sich die Einschätzung der Einstellungspolarität nach der Semantik des Ursprungslexems richtet, wenn der Intensivierer ein semantisch leeres Pseudoadjektiv modifiziert.


%In einer weiteren Untersuchung wurde der Annahme des wahrnehmbaren Unterschiedes zwischen deskriptiven und expressiven Intensivierern noch einmal gesondert auf den Grund gegangen. Das in \mbox{Kapitel \ref{Experimente}} dargelegte Experiment fußte auf den folgenden beiden Hypothesen:
%
%\begin{itemize}
%\item[\textbf{\textsc{h4}}] Expressive Intensivierer werden mit einem höheren Grad an Skalarität
%assoziiert als deskriptive Intensivierer. 
%
%\item[\textbf{\textsc{h5}}] Expressive Intensivierer geben eine Einstellungsbekundung des Sprechers
%zu dem im Satz beschriebenen Sachverhalt wieder und werden in der Folge mit einem höheren Grad an Sprechereinstellung assoziiert als deskriptive Intensivierer. 
%\end{itemize}
%
%\noindent Der Hypothese H4 nach geben expressive Intensivierer zum einen einen höheren Grad einer Merkmalsausprägung an und werden daher in Bezug auf Skalarität auf einer Skala höher verortet als deskriptive Intensivierer (vgl. \citealt[22]{Paradis.1997}; \mbox{\citealt[133]{Gutzmann.2019a}}). Zum anderen dienen sie H5 zufolge ferner dem Ausdruck der persönlichen Einstellung des Sprechers, weswegen ihnen auch im Hinblick auf dieses Merkmal vermutlich eine höhere Skalenposition zugewiesen wird (vgl. \citealt[315]{Lang.1983}; \citealt[57]{Ortner.2014}).


%Die Gültigkeit dieser Hypothesen wurde in zwei Rating-Experimenten mit parallelem Aufbau getrennt voneinander untersucht. Dabei beurteilte eine Personenstichprobe Sätze mit deskriptiven und expressiven Intensivausdrücken hinsichtlich Skalarität (\mbox{Experiment 2a}) und Sprechereinstellung (Experiment 2b). Durch die Parallelität konnten die Ergebnisse der beiden Experimente miteinander verglichen und auf eine mögliche Korrelation zwischen Skalarität und Sprechereinstellung hin überprüft werden.


%Die Ergebnisse der Experimente bekräftigten die beiden zu testenden \mbox{Hypothesen H4 und H5}, wurden Sätze mit expressiven Intensivierern doch sowohl mit Blick auf Skalarität als auch auf Sprechereinstellung von den Versuchspersonen auf der Skala tendenziell höher positioniert als Sätze mit deskriptiven Intensivierern. Wenngleich der Unterschied bei der in \mbox{Experiment 2a} abgefragten Skalarität mit ca. 4,50\,\% nicht übermäßig stark ausgeprägt war (vgl. \mbox{Tabelle \ref{tab:Skalenwert1}}), war er bei der in Experiment 2b fokussierten Sprechereinstellung mit ca. 11,30\,\% dagegen umso größer (vgl. \mbox{Tabelle \ref{SkaWe_E2}}). Bei der anschließend vorgenommenen Einzelbetrachtung der Intensivierer fiel im Kontext von Experiment 2a eine vorherrschende Trennung nach deskriptiv und expressiv auf (vgl. \mbox{Tabelle \ref{Skalenwert3}}). Dennoch war in beiden Intensivierklassen eine Ausnahme auszumachen, nämlich \textit{überaus} aufseiten der deskriptiven sowie \textit{voll} aufseiten der expressiven Intensivierer (vgl. \mbox{Abbildung \ref{Bar1}}). Während ersterem mit durchschnittlich 62,36\,\% ein verhältnismäßig hoher Skalaritätsgrad zugesprochen wurde, wurde \textit{voll} mit durchschnittlich 59,69\,\% von den Versuchspersonen dahingehend auf der Skala merklich niedrig verortet. Für diesen Befund wurden zwei Interpretationsmöglichkeiten der Hypothese H4 offengelegt, von denen die gemäßigtere eine mögliche Erklärung für dieses scheinbar abweichende Verhalten der beiden Ausdrücke bot (vgl. \mbox{Tabelle \ref{schwach}}). Die Einzelbetrachtung der Intensivierer von Experiment 2b zeigte hingegen keine solchen vermeintlichen Ausreißer, was bedeutet, dass sich in Bezug auf Sprechereinstellung erwartungsgemäß eine markante Trennung nach Intensiviererklasse herauskristallisierte (Tabelle \ref{Einzeln_2} und \mbox{Abbildung \ref{Bar2}}).


%Um zu überprüfen, ob die beobachteten prozentualen Unterschiede von statistischer Relevanz waren, wurden die Intensivierer mittels \textit{t}-Tests paarweise miteinander verglichen. Im Hinblick auf die in Experiment 2a erfragte Skalarität hoben sich insbesondere drei Intensivierer stark von den anderen ab (vgl. Tabelle \ref{Matrix_1}): der deskriptive Intensivierer \textit{ziemlich} sowie die expressiven Intensivierer \textit{mega} und \textit{voll}.  Dadurch, dass ersterer \textit{qua} Semantik kein Höchstmaß der bezeichneten Eigenschaft, sondern einen eher moderaten Intensivierungsgrad ausdrückt, schienen dessen signifikanten Unterschiede beim Vergleich mit den anderen Ausdrücken kaum verwunderlich. Die signifikanten Unterschiede von \textit{mega} wurden auf die anhaltende Aktualität des Ausdrucks zurückgeführt, die sich offenbar in einem als hoch empfundenen Skalaritätsgrad niedergeschlagen hat. Dagegen wurden die signifikanten Unterschiede von \textit{voll} als Folge von pragmatischer Abnutzung durch persistenten Gebrauch gesehen, die dessen vergleichsweise niedrigen Skalenwert verantwortet zu haben schien. Noch ausgeprägter waren die Ergebnisse der Paarvergleiche der Intensivierer zu der in \mbox{Experiment 2b} untersuchten Sprechereinstellung, unterschieden sich die beiden Intensiviererklassen hierbei doch nahezu komplett voneinander (vgl. \mbox{Tabelle \ref{Sign2}}). Lediglich der deskriptive Intensivierer \textit{sehr} stellte an dieser Stelle eine Ausnahme dar, die durch die ungewöhnlich hohe Skalenposition des Ausdrucks bedingt war. Dieser wurde im Rahmen der Ergebnisinterpretation auf dessen hohen Lexikalisierungsgrad zurückgeführt, durch den \textit{sehr} nicht nur der klassische deskriptive Intensivierer des Deutschen ist, sondern den Ergebnissen der Untersuchung nach auch mit einem hohen Grad an Sprechereinstellung assoziiert zu werden scheint. Das Regressionsmodell, das für die in \mbox{Experiment 2a} erfragte Skalarität berechnet wurde, enthielt signifikante Haupteffekte für die Prädiktoren \textsc{Intensiviererklasse} und \textsc{Position} sowie deren Interaktion (vgl. \mbox{Tabelle \ref{LinReg1}}). Somit war der Unterschied, der sich zwischen den beiden Intensiviererklassen abzeichnete, signifikant (\mbox{$\chi$$^{2}$\,(1) = 8,81}, \mbox{\textit{p} < 0,01)}. Zudem wurde ein Item mit fortschreitender Position im Experiment auf der Skala signifikant niedriger verortet (\mbox{$\chi$$^{2}$\,(1) = 16,20}, \mbox{\textit{p} < 0,001)}, während die signifikante Zwei-Wege-Interaktion weiterhin angab, dass sich die eingeschätzten Skalenwerte der Intensiviererklassen über das Experiment hinweg sukzessive annäherten (\mbox{$\chi$$^{2}$\,(1) = 4,17}, \mbox{\textit{p} < 0,05)}. Dies wurde als ein Indiz für einen Gewöhnungseffekt gesehen, der sich dem Anschein nach bei den Versuchspersonen infolge der uniformen Aufgabenstellung eingestellt hat. Auch in dem für Experiment 2b berechneten finalen Modell war ein signifikanter Haupteffekt für den Prädiktor \textsc{Intensiviererklasse} auszumachen (vgl. Tabelle \ref{LinReg2}). Demnach war der verzeichnete Unterschied zwischen deskriptiven und expressiven Intensivierern, der sich in Bezug auf Sprechereinstellung zeigte, statistisch signifikant (\mbox{$\chi$$^{2}$\,(1) = 14,51}, \mbox{\textit{p} < 0,001)}. Kurz gefasst sprachen sich die Ergebnisse der Experimente für die Annahmen von \citet[22]{Paradis.1997} und \mbox{\citet[133]{Gutzmann.2019a}} sowie \mbox{\citet[315]{Lang.1983}} und \mbox{\citet[57]{Ortner.2014}} und folglich die Gültigkeit der zu testenden Hypothesen H4 und H5 aus.


%Ergänzend wurde überprüft, ob die beiden Variablen Skalarität und Sprechereinstellung korreliert waren. Während der \mbox{Korrelationskoeffizient \textit{r}} schon auf Unkorreliertheit hindeutete (|r| = 0,24), wurde dieses Ergebnis durch einen nicht signifikanten \textit{t}-Test untermauert (\mbox{\textit{t}\,(6) = 0,59}, \mbox{\textit{p} > 0,05}). Demzufolge waren Skalarität und Sprechereinstellung voneinander unabhängig. Für diesen Befund wurden unterschiedliche Interpretationsmöglichkeiten geliefert, von denen die erste auf ein methodisches Defizit zurückgeführt wurde: So wurde in der Aufgabenstellung von Experiment 2b nicht nach dem Grad an Sprechereinstellung, sondern dem Grad an Sprecherbewertung gefragt, was an der Stelle zu eventuell verzerrten Ergebnissen geführt hat. Dass die Untersuchungsmethode mit dahingehend optimierter Fragestellung zu den gewünschten Ergebnissen führt, konnte durch das in \mbox{Kapitel \ref{Experiment0}} präsentierte Experiment gezeigt werden. Daneben wurde die Möglichkeit angeführt, dass die Versuchspersonen den Skalaritätsgrad der Intensivierer fälschlicherweise als Sprechereinstellung interpretiert und den deskriptiven Ausdrücken in der Konsequenz eine verhältnismäßig hohe Skalenposition zugewiesen haben. Eine weitere Interpretationsmöglichkeit bezog sich dagegen auf die vorgenommene Klasseneinteilung der Intensivierer. Vor diesem Hintergrund wurde spekuliert, dass die Experimente mit anderen, (noch) aktuelleren Ausdrücken zu anderen Resultaten und möglicherweise zu einer signifikanten Korrelation zwischen den Variablen Skalarität und Sprechereinstellung geführt hätten. Da sich diese Vermutung im Nachgang an die Untersuchung jedoch nicht zweifelsfrei einschätzen lässt, muss deren Überprüfung nachfolgenden Experimenten vorbehalten bleiben.


%Zusätzlich wurde mithilfe weiterer \textit{t}-Tests untersucht, ob zwischen den in den Experimenten verwendeten Pseudowörtern (vgl. Kapitel \ref{Experiment0}) und realen Wörtern (vgl. Kapitel \ref{Experimente}) ein korrelativer Zusammenhang bestand. Zur Überprüfung wurden die Mittelwerte der Schnittmenge der Wörter für Skalarität und Sprechereinstellung miteinander verglichen. Während die berechneten \textit{t}-Tests zunächst für beide signifikant waren, wurden die Werte des deskriptiven Intensivierers \textit{sehr} zur Kontrolle ausgeklammert und nur noch die Mittelwerte der expressiven Intensivierer dahingehend getestet. Die Ergebnisse dieser Tests waren indes nicht signifikant: Skalarität (\mbox{\textit{t}\,(3) = 28,12}, \mbox{\textit{p} > 0,05}) und Sprechereinstellung (\textit{t}\,(3) = 15,54, \mbox{\textit{p} > 0,05}). Demzufolge waren die zuvor festgestellten signifikanten Unterschiede alleine durch die Mittelwerte von \textit{sehr} bedingt. Dagegen spielte es für die expressiven Intensivierer keine Rolle, ob sie semantisch leere Pseudoadjektive oder existente Adjektive modifizierten -- sie verhielten sich in beiden Kontexten parallel. Aufgrund der nachgewiesenen Parallelität konnte angenommen werden, dass die Semantik der realen Wörter die Skalenpositionierung bei den in Kapitel \ref{Experimente} geschilderten Experimenten nicht unerwünscht beeinflusst hat und die Ergebnisse beider Untersuchungen gleichermaßen repräsentativ sind. Stattdessen schien die Einschätzung der Items auf der Skala, wie intendiert, ausschließlich durch die Intensivierer zustande gekommen zu sein, wodurch sich für anschließende Untersuchungen die Wahlmöglichkeit des lexikalischen Materials ergibt.


%Zusammengefasst belegten die in der Arbeit vorgestellten Ergebnisse der empirischen Untersuchungen nahezu aller zur Diskussion stehenden Hypothesen: Während die Hypothese H1-Freq nach \citet[2]{Hauschild.1899} sowohl durch einen signifikanten Unterschied zwischen den Frequenzgruppen als auch eine negative Korrelation zwischen Frequenz und Sprechereinstellung gestützt wurde, legten die im Korpus rezent aufgetretenen Intensivierer \textit{hammer}, \textit{krass} und \textit{mega} das Zutreffen der Hypothese H2 nach \mbox{\citet[274]{Baumgarten.1908}} nahe, nach der infolge des Expressivitätsverlusts alter Intensivierer permanent neue Ausdrücke erforderlich sind. Abgesehen davon konnte auch die Hypothese H3 für gültig befunden werden, die sich aus einer Behauptung von \citet[15f.]{Nouwen.2020} heraus ableitete, der beschreibt, dass die Semantik vieler Intensivierer gänzlich ausgebleicht ist. Wie die Ergebnisse des in Kapitel \ref{Experiment0} vorgestellten Experiments zeigten, ist dies im Falle der Modifikation eines Pseudowortes nicht der Fall. So wurde die Einschätzung der Einstellungspolarität in Übereinstimmung mit H3 in der Tat durch die Semantik des Ursprungslexems hervorgerufen. Darüber hinaus fanden sich ebenso klare Hinweise für das Zutreffen der \mbox{Hypothesen H4 und H5} nach \citet[22]{Paradis.1997} und \mbox{\citet[133]{Gutzmann.2019a}} sowie \citet[315]{Lang.1983} und \citet[57]{Ortner.2014}, nach denen expressive Intensivierer gleichermaßen mit einem höheren Grad an Skalarität und Sprechereinstellung einhergehen. Folglich konnten lediglich für die \mbox{Hypothese H1-Temp} nach \citet[59]{Tobler.1868}, nach der langer Gebrauch zulasten der Expressivität von Ausdrücken geht, keine harten empirischen Evidenzen gefunden werden. Da die Ergebnisse der in Kapitel \ref{Experiment0} dargelegten Untersuchung keine signifikanten Unterschiede zwischen den Persistenzgruppen in Bezug auf Sprechereinstellung zutage brachten, musste davon ausgegangen werden, dass die im Korpus neu aufgetretenen Intensivierer \textit{hammer}, \textit{krass} und \textit{mega} von den Versuchspersonen als nicht signifikant expressiver empfunden wurden als die älteren Intensivierer (n = 13). Dieses Resultat ist in Anbetracht der niedrigen Anzahl an neuen Intensivierern allerdings nur eingeschränkt aussagekräftig, hat das daraus resultierende Ungleichgewicht im Umfang der Persistenzgruppen doch vermutlich auf die inferenzstatistischen Berechnungen ein- und dem erwarteten Effekt entgegengewirkt. Demnach ist vorstellbar, dass die Ergebnisse mit einem ausgeglichenen Verhältnis aus alten und neuen Intensivierern an dieser Stelle anders ausgesehen hätten. Diese Überlegung könnte gegebenenfalls in Folgeexperimenten \mbox{berücksichtigt werden}.


%Abschließend gilt es, die der Arbeit übergeordnete, titelgebende Frage zu beantworten, die auf eine Behauptung von \citet[133]{Gutzmann.2019a} zurückgeht: Ist \textit{sau cool} cooler als \textit{sehr cool}? Wie die Ergebnisse der experimentellen Untersuchungen mit deutlicher statistischer Signifikanz gezeigt haben, scheint dies in der Tat zutreffend zu sein. In diesem Sinne kann  schlussendlich ganz klar festgehalten werden: \textit{Sau cool} ist zweifellos cooler als \textit{sehr cool}.





%empirische Datengrundlage: Ergebnisse der Korpusstudie aus Kapitel 3, bei der es sich um eine großangelegte diachrone Untersuchung von Intensivierern handelt, die im Folgenden (Kapitel 5 und 6) experimentell untersucht wurden
%korpusgestützt
%breite empirische Grundlage


%Es bleibt jedoch (empirisch) zu klären,...
% Es scheint die Annahme plausibel, dass...
%Dieses Kapitel dient der Diskussion der Annahmen, die im Hinblick auf Intensivierung als einschlägig aufgeworfen wurden. 

%Ausgehend von den experimentellen Ergebnissen
%Zentralpunkt
%von Belang sein

%Ob sich dieser Befund durch weitere und größer angelegte Untersuchungen erhärten lässt, muss an dieser Stelle offenbleiben.

%Damit wird durch die Daten die eingangs formulierte Hypothese nicht gestützt
%Propositionale (Sprecher-)Einstellung, d. h. dem Sprecher werden gewisse Gefühlsregungen und Einstellungen zugeschrieben
%Deskriptive und expressive Adjektive konstrastiert wurden

%Wolski (1989: 349): „Als Bedeutung eines sprachlichen Ausdrucks werden die Regeln für den Gebrauch, seine semantischen Gebrauchsregeln aufgefaßt.“

%Möglicher Expressivitätsverlust bei derb


%Die wichtigsten Ergebnisse / Erkenntnisse dieser Arbeit beziehen sich auf


%Berz (1953: 13): "So wie hier \textit{furchtbar} [in \textit{furchtbar nett}, Anmerkung J. S.] seinen eigentlichen Sinn verloren hat und nur noch zur Steigerung des zweiten Begriffs dient, so ist das erste Glied beim T. st. [= Typus \textit{steinreich}, Anmerkung J. S.] seines Bedeutungsinhaltes entleert und dient zur Verstärkung des Grundwortes." --> Inwiefern die ursprünglich negative Bedeutung noch wahrgenommen wird, bedarf weiterer Überprüfung (vgl. Ergebnisse E-Mail-Umfrage, die auf Gegenteiliges hindeuten)

% Os (1989: 190): "Auch Biedermann (1969) gibt trotz seiner Subklassifizierung nach Motivationssemen keine befriedigende Antwort auf die Frage, inwieweit die lexikalische Bedeutung der einzelnen IM (= Intensivierunsmittel) in den Hintergrund gedrängt ist. Ein Maß dafür wäre die Weite der Kontextverträglichkeit."



%als Ende der Arbeit Becher 1907: 273: "Schließlich ist nur noch darauf hinzuweisen, was uns zum Gebrauch der steigenden Ausdrücke veranlassen kann und auf welchen Gebieten sie besonders üblich sind. Der Grund ist immer derselbe, das Bedürfnis der Hervorhebung, der Verstärkung. Dazu kann uns veranlassen: 1. das Bedürfnis, den Ausdruck herzlich zu gestalten; 2. das Bedürfnis der Rhetorik [...]; 3. der sogenannte Kurialstil und endlich 4. die Reklame." --> Auch wenn seit diesem Zitat mehr als ein Jahrhundert vergangen ist, hat sich am Inhalt jedoch nichts verändert.
